O presente trabalho estabelece uma arquitetura de segurança integrada para 
redes LoRaWAN, combinando criptografia leve por meio do padrão Ascon e 
detecção de anomalias via modelos de regressão temporal. Embora os objetivos 
definidos para este projeto tenham sido tratados nas etapas descritas nos 
capítulos anteriores, dentro das limitações ali declaradas, a natureza do 
problema abre caminhos relevantes para investigações futuras.

A perspectiva de maior interesse imediato diz respeito à portabilidade da 
fase de inferência dos modelos treinados para o próprio gateway da rede. 
Neste trabalho, o treinamento e a avaliação dos modelos são realizados em 
ambiente computacional convencional, decisão justificada pelas restrições de 
hardware incompatíveis com o processo de treinamento. Contudo, uma vez 
obtido um modelo com desempenho satisfatório, torna-se pertinente investigar 
se sua fase de inferência pode ser executada diretamente no gateway, 
eliminando a necessidade de transmissão dos dados brutos para uma camada 
externa de processamento.

Essa investigação demandaria a análise da complexidade do modelo resultante 
— em termos de número de parâmetros, tamanho serializado e tempo de 
inferência — em relação às restrições de memória RAM e armazenamento em 
\textit{flash} do hardware adotado. Ferramentas como o TensorFlow Lite for 
Microcontrollers e a biblioteca \textit{emlearn} permitem a conversão de 
modelos treinados em Python para código C otimizado, compatível com 
microcontroladores como o ESP32 \cite{ray2022}. Técnicas de compressão como 
quantização de pesos e poda de conexões poderiam ser aplicadas para reduzir 
o custo computacional sem perda significativa de desempenho \cite{jacob2018}.

Uma segunda perspectiva relevante envolve a extensão da abordagem de 
detecção para incorporar dados rotulados de ataque, como os disponíveis no 
\textit{dataset} GUARD. Com resultados de regressão consolidados sobre o 
Tour Perret, seria possível investigar em que medida os resíduos produzidos 
pelos modelos se comportam de forma diferenciada diante de \textit{frames} 
reconhecidamente manipulados, abrindo caminho para uma avaliação mais 
rigorosa da capacidade de detecção em cenários com ataques conhecidos.

Por fim, a arquitetura proposta poderia ser estendida para contemplar 
mecanismos de atualização contínua do modelo, permitindo que o sistema se 
adapte a mudanças graduais no comportamento da rede — como variações 
sazonais de sinal ou substituição de dispositivos — sem comprometer a 
sensibilidade à detecção de anomalias. Essa capacidade de adaptação é 
especialmente relevante em implantações de longa duração, características 
de aplicações industriais e de cidades inteligentes.
